\chapter*{Anhang zur Dokumentation der verwendeten KI-basierten Hilfsmittel}
\vspace{-6em}{\Huge\color{dkfzgray2} Appendix documenting the AI-based tools used}
\addcontentsline{toc}{chapter}{Anhang zur Dokumentation der verwendeten KI-basierten Hilfsmittel}
\par\vspace{2em}

% ---------------------------------------------------------------------------
% ANLEITUNG (vor Einreichung löschen):
%
% Aufbau nach Muster 11 b der Abfassungsrichtlinien (Version 2.5). Dieser Anhang
% steht direkt im Anschluss an die Erklärung (11_Angabe_KI.tex). Für JEDES
% verwendete KI-Tool (z. B. Claude Code, ChatGPT, DeepL, Elicit, GitHub Copilot,
% ...) eine eigene Tabellenzeile anlegen:
%   KI-Tool                    Name des Werkzeugs
%   Version                    Modell/Version oder Nutzungszeitraum (z. B. 10.2026)
%   Verwendete Funktionalität  z. B. LLM, systematische Literaturrecherche
%   Nutzungszweck              z. B. Textkorrektur, Ergänzung der manuellen
%                              Literaturrecherche, Konvertierung eines Entwurfs
%   Einsatzbereich             konkrete Kapitel/Abschnitte, nicht "überall"
%   Art der Validierung        wie die Ergebnisse geprüft wurden, z. B.
%                              Diff-Ansicht, manuelle Quellenprüfung in PubMed,
%                              gezielte manuelle Überprüfung, eigene Texterstellung
%
% KI-Anwendungen, die bereits in "Material und Methoden" ausführlich
% beschrieben sind (z. B. KI als Forschungsmethode), brauchen hier nur einen
% Verweis auf das jeweilige Unterkapitel.
%
% Die Zeile zu Claude Code ist bereits vorbereitet, weil diese Vorlage für die
% Arbeit mit Claude Code gedacht ist. Der Absatz unter der Tabelle zur
% Arbeitsumgebung (heidelberg-med-thesis-claude-template) und zum zugrunde
% liegenden LaTeX-Template von Ayala, Kausch, Klein und Thomas ist die nach
% CC BY 4.0 erforderliche Namensnennung und muss erhalten bleiben.
% ---------------------------------------------------------------------------

{\footnotesize
\renewcommand{\arraystretch}{1.3}
\noindent\begin{tabularx}{\textwidth}{@{}*{6}{>{\raggedright\arraybackslash}X}@{}}
\toprule
\textbf{KI-Tool} & \textbf{Version} & \textbf{Verwendete Funktiona\-lität} & \textbf{Nutzungs\-zweck} & \textbf{Einsatz\-bereich} & \textbf{Art der Validierung} \\
\midrule
Claude Code (Anthropic)
  & TODO: Modell und Nutzungs\-zeitraum
  & TODO (z.\,B. LLM-Agent mit Zugriff auf die LaTeX-Projekt\-dateien)
  & TODO (z.\,B. Konvertierung des Entwurfs in die LaTeX-Vorlage, Textkorrektur)
  & TODO (konkrete Kapitel/\-Abschnitte)
  & TODO (z.\,B. Diff-Ansicht jeder Änderung, gezielte manuelle Überprüfung) \\
\midrule
TODO: weiteres Tool & & & & & \\
\bottomrule
\end{tabularx}
}

\bigskip

\noindent Claude Code wurde innerhalb der frei verfügbaren Arbeitsumgebung
„heidelberg-med-thesis-claude-template“ von Peter Maximilian Full
eingesetzt (\url{https://github.com/pmfull/heidelberg-med-thesis-claude-template},
Lizenz CC BY 4.0). Diese stellt die LaTeX-Vorlage dieser Arbeit sowie
projektspezifische Arbeitsanweisungen für Claude Code bereit, die u.\,a. die
Einhaltung der Abfassungsrichtlinien der Medizinischen Fakultät Heidelberg und
der KI-Leitlinie der Universität Heidelberg vorgeben und das Erfinden von
Quellen untersagen, sowie einen Prüf-Agenten, der die Arbeit auf Formvorgaben
und Kohärenz prüft, ohne sie zu verändern. Die LaTeX-Vorlage basiert auf dem
LaTeX-Template von Leonardo Ayala, Lisa Kausch, André Klein und Sarina Thomas
(\url{https://www.overleaf.com/latex/templates/heidelberg-university-medical-faculty-doctoral-thesis/fzxdxksbfqvm},
Lizenz CC BY 4.0).

\emph{\small\color{dkfzgray2} Claude Code was used within the freely available
working environment “heidelberg-med-thesis-claude-template” by Peter
Maximilian Full
(\url{https://github.com/pmfull/heidelberg-med-thesis-claude-template},
licensed under CC BY 4.0). It provides the LaTeX template of this thesis as
well as project-specific instructions for Claude Code that, among other things,
require compliance with the writing guidelines of the Medical Faculty
Heidelberg and the AI guideline of Heidelberg University and prohibit
fabricating sources, and a review agent that checks the thesis for formal
requirements and coherence without modifying it. The LaTeX template is based on
the LaTeX template by Leonardo Ayala, Lisa Kausch, André Klein, and Sarina
Thomas
(\url{https://www.overleaf.com/latex/templates/heidelberg-university-medical-faculty-doctoral-thesis/fzxdxksbfqvm},
licensed under CC BY 4.0).}
